% ============================================================
% CHAPTER 2: BACKGROUND
% All citations in this chapter reuse entries already present in
% bibliography.bib and verified during the Chapter 3 literature
% review. No new references are introduced here.
% ============================================================

Multi-agent debate is a multi-agent system, and its documented failures are
dynamical: they are properties of how positions evolve and influence one another
across rounds, not of any single agent's competence at a single moment. Naming
those failures precisely, and measuring them, requires a vocabulary for
dynamical systems. Most readers of this thesis are practitioners who know what
an LLM agent is and what multi-agent debate does; they are far less likely to
know what an attractor is, what a basin of attraction selects, or how influence
concentrates in a network through purely structural means. This chapter
therefore spends most of its space on the unfamiliar ground.
\autoref{sec:bg:complex} introduces complex systems and self-organization from
first principles: emergence, attractors, bistability, and positive feedback.
\autoref{sec:bg:motifs} casts those ideas as formal motifs and introduces
opinion dynamics, the branch of complex-systems theory that studies how
influence propagates and concentrates among interacting agents, and which
provides the closest formal model of what debate agents do.
\autoref{sec:bg:mad} fixes the familiar parts: what an LLM agent means in the
sense used here, how multi-agent debate works, and which failure modes motivate
this study. Formal definitions of the four detectors are deferred to
\autoref{sec:method}; the present chapter only establishes the ideas they rest
on.

% ============================================================
\subsection{Complex Systems and Self-Organization}
\label{sec:bg:complex}
% ============================================================

A \emph{complex system} is a collection of many interacting components that
collectively exhibit behaviour which cannot be read off from any individual
component in isolation. The behaviour is called \emph{emergent}: it is a
property of the interaction structure, not of any single unit. When such
emergent behaviour is ordered and coordinated and arises without any central
authority imposing it, the system is called \emph{self-organizing}
\citep{ashby1947principles}. No component decides the global outcome, yet a
global outcome reliably appears. A flock of birds, a developing embryo, and a
financial market are all self-organizing: each component follows local rules
about its immediate neighbours, yet the collective as a whole exhibits
coherent global structure that no individual plans or directs.

The natural language for self-organizing systems is that of dynamical systems
theory. The state of a system at a given moment is a point in a space whose
dimensions correspond to all quantities that describe the system; as time
passes, the system traces a \emph{trajectory} through that space. Not all
parts of the space are equally visited: over time, trajectories tend to settle
onto sets called \emph{attractors} \citep{strogatz2015nonlinear}. A
\emph{fixed-point attractor} is a state the system reaches and stays in: a
consensus, a locked configuration. A \emph{limit cycle} is a closed orbit the
system revisits indefinitely without settling, a permanent oscillation in which
the same sequence of states repeats in a regular cycle. The region of initial
conditions that eventually leads to a given attractor is its \emph{basin of
attraction}. Two initial conditions in the same basin converge to the same
long-run behaviour; two initial conditions in different basins do not.

A system is \emph{bistable} (or more generally \emph{multistable}) if two or
more distinct attractors coexist for the same parameter values
\citep{feudel2008complex}. Which attractor is reached is then determined
entirely by the initial condition: two runs of the same system with different
starting points can land on different long-run outcomes for purely dynamical
reasons, not because anything about the system itself changed. The boundary
between basins, the \emph{separatrix}, need not lie at an obvious midpoint. In
social systems, the corresponding phenomenon is a tipping point: a critical
minority is sufficient to move the collective from one stable state to another,
well below the majority that naive equilibrium reasoning would require.
\citet{centola2018tipping} establish experimentally that this critical mass is
roughly 25\% of the population, that it is robust across norm types, and that
longer interaction memories raise the threshold.

Two further regimes are central to this thesis. \emph{Self-reinforcement}, or
positive feedback, is the tendency of a state to amplify its own persistence:
an early lead grows because it is a lead, independently of whether it is
intrinsically superior. The canonical example is technology lock-in through
increasing returns, where a product's market share generates adoption that
further increases that share \citep{arthur1989competing}. Positive feedback
does not select the best option; it selects the option that accumulated an
early advantage. The complementary phenomenon is the \emph{concentration of
influence} onto a single component. In a network of interacting agents, one
node may come to disproportionately determine the collective state, not because
it is better informed, but because the interaction structure channels influence
through it. Both regimes, positive feedback and influence concentration, are
detectable from ordinary time-series records of the system's state, which is
why they admit an empirical instrument.

% ============================================================
\subsection{Dynamical Motifs and Opinion Dynamics}
\label{sec:bg:motifs}
% ============================================================

The patterns described above can be made precise as \emph{motifs}: recurring,
detectable signatures of a system's organization. The term carries two
complementary meanings that this thesis uses carefully. A \textbf{structural
motif} \citep{milo2002network, alon2007network} is a recurring pattern in the
\emph{wiring}: a small subgraph of who-connects-to-whom that appears in a
network more often than chance and generates characteristic input--output
behaviour. The archetypal example is the feed-forward loop, a three-node
circuit whose sign configuration makes it a sign detector or a pulse filter.
A \textbf{dynamical motif} \citep{driscoll2024dynamical} is instead a
recurring pattern in the \emph{trajectory}: an attractor, a periodic orbit, a
separatrix between basins of attraction. It is a statement about how the
system behaves over time rather than how it is wired. The two are linked,
since structural motifs are studied precisely because each generates a
characteristic dynamical regime, and they map onto the two levels at which
phenomena live in this thesis, the agent level and the system level.

This thesis detects four motifs, each selected because it has a precise,
testable signature and corresponds to a documented failure mode of multi-agent
debate. \emph{Self-reinforcement} is positive feedback at the agent level: an
agent that holds a position through a round emerges with higher rather than
lower confidence, a positive-autoregulation structural motif whose dynamical
signature is a rising confidence slope within vote-stable runs.
\emph{Limit cycles} are self-sustained periodic oscillations: a vote sequence
that revisits the same states in a repeating order without converging, a
dynamical motif readable from the discrete vote channel.
\emph{Bistability} is the coexistence of multiple stable consensus outcomes
across repetitions of the same task, a dynamical motif whose signature is a
multimodal outcome distribution that persists as the initial vote composition
varies. \emph{Dominance} is the concentration of causal influence onto a
single agent, reflected in a heavy-tailed distribution over an emergent
influence graph, a structural motif in the influence network that the debate
generates. Together the four span both senses of motif and both levels of
description.

The propagation and concentration of influence that underlie self-reinforcement
and dominance have a formal home in \emph{opinion dynamics}, which studies how
agreement and social power evolve in a network of interacting agents. The
classical model due to \citet{degroot1974consensus} treats each agent's
opinion as a weighted average of its neighbours' opinions at the previous step.
Under mild connectivity conditions this process reaches consensus: all agents
converge to the same value, which is a weighted combination of their initial
opinions. The coefficients in this combination are called \emph{social power}
\citep{jia2015opinion}, and they are not generally equal: an agent that is
upstream of many others accumulates disproportionate weight in the final belief,
even if its initial opinion is no more accurate than its peers'. This is the
formal analogue of the dominance motif: influence concentration arising from
network position rather than from informational superiority.

Our agents differ from the classical opinion dynamics setting in one important
respect: they hold a latent continuous confidence, elicited as an integer from
1 to 10, but cast a discrete vote. This \emph{continuous-opinion, discrete-action}
structure is a studied variant of the classical framework
\citep{martins2008coda}: social influence propagates through the continuous
channel while discrete commitment is cast at decision time. The separation is
analytically useful. Self-reinforcement is detectable on the continuous channel,
where a slope is well-defined; cycles and attractor statistics are detectable
on the discrete channel, where the state space is small enough to characterise
completely. A group of homogeneous LLM agents debating over multiple rounds
occupies exactly this regime: components following local interaction rules and
generating global collective behaviour, the defining signature of a
self-organizing complex system.

% ============================================================
\subsection{LLM Agents and Multi-Agent Debate}
\label{sec:bg:mad}
% ============================================================

A large language model (LLM) is a neural network, in current practice a
Transformer \citep{vaswani2017attention}, trained to predict the next token in
a sequence and subsequently aligned to follow instructions through supervised
fine-tuning and reinforcement learning from human feedback
\citep{ouyang2022instructgpt}. An \emph{LLM agent}, in the sense used
throughout this thesis, is an LLM placed in a loop: it receives a state
comprising the task, its own prior position, and messages from other agents; it
produces an action consisting of a written argument and a vote; and the
resulting state is fed back on the next step. Sampling is stochastic: at a
positive temperature the model does not return a single deterministic answer
but draws from its output distribution, so two agents given the identical
prompt can produce different responses. This stochasticity is the source of
initial diversity in a debate and, as later chapters show, the raw material for
a bistability analysis. Because we use a single model family throughout, all
four agents share the same priors and the same competence; any diversity in
the group is a product of sampling, not of heterogeneous models.

\emph{Multi-agent debate} (MAD) is a protocol in which several LLM agents
independently propose an answer to a question and then exchange and revise
their reasoning over multiple rounds before the group commits to a final
decision \citep{du2024improving}. The motivating intuition is an
error-correction argument: if agents make partially independent mistakes,
deliberation gives the group a chance to surface and correct them, much as an
ensemble of independent voters can outperform its average member
\citep{grofman1983thirteen}. Reported gains on factual and mathematical
reasoning \citep{du2024improving, liang2024encouraging, chen2024reconcile} and
the use of debate panels as more reliable evaluators \citep{chan2024chateval}
follow from this idea. Debate designs vary along three structural axes that
recur throughout this thesis. The \emph{communication topology} fixes who can
read whom: a fully-connected group in which every agent sees every other versus
sparser arrangements such as a star, where peripheral agents communicate only
through a central hub. Sparser topologies can match fully-connected debate at
lower token cost \citep{li2024sparsetopology}, so topology is a genuine design
decision. The \emph{memory window} fixes how much interaction history each
agent retains from one round to the next. The \emph{decision protocol} fixes
how individual positions are aggregated into a group answer; this choice alone
can change which task types benefit from debate \citep{kaesberg2025voting}. We
treat topology and memory window as controlled experimental axes and hold the
protocol fixed.

Crucially for this thesis, the same interaction that enables error correction
also admits systematic failure. A purely cooperative framing can collapse the
debate into mutual agreement that entrenches an initial mistake, the
\emph{Degeneration-of-Thought} failure mode \citep{liang2024encouraging}.
When two models debate, both tend to grow more confident that they are winning
as the exchange proceeds, with confidence escalating away from a rational
baseline even absent new evidence \citep{prasad2025certain}, and agents
frequently abandon a correct answer to conform to a confident peer
\citep{wynn2025talk}. This behaviour has been formalised as inter-agent
sycophancy \citep{yao2025sycophancy}, the multi-agent extension of the
single-agent sycophancy induced by preference-based alignment
\citep{sharma2023sycophancy}. When the information needed for the correct
answer is distributed across agents rather than shared, groups lock onto the
shared evidence and fail to pool the private pieces, so collective accuracy
collapses far below what a single fully-informed agent achieves
\citep{li2026hiddenbench}; this is the LLM analogue of the hidden-profile
effect long documented in human groups \citep{stasser1985pooling}. These
pathologies are not accidents of a particular prompt. They are properties of
how positions evolve and influence one another over rounds, and the vocabulary
of the preceding two sections is now precise enough to name them: escalating
confidence is a self-reinforcement signature, premature consensus is a
fixed-point attractor reached through basin selection, persistent disagreement
resolving only through capitulation is a dominance pattern, and unresolved
cycling is a limit-cycle regime.
